\begin{table}[H]
\color{black}
\captionsetup{labelfont={color=black},textfont={color=black}}
\caption{Scale of the Response Discrepancy Across Signal Durations}
\label{tab:app-response-scale}
\centering\footnotesize
\setlength{\tabcolsep}{2pt}
\renewcommand{\arraystretch}{1.12}
\begin{tabular*}{\columnwidth}{@{\extracolsep{\fill}}lrrrrr@{}}
\toprule
Date & \shortstack{Duration\\(h)} & \shortstack{$E_0$\\(MWh)} & \shortstack{$\delta P_C$\\(kW)} & \shortstack{$\epsilon_C$\\(\%)} & \shortstack{$\epsilon_A$\\(\%)}\\
\midrule
07-27 & 1 & 3.999 & 0.000 & 0.000 & -0.001\\
 & 2 & 8.077 & -0.690 & -0.017 & -0.027\\
 & 4 & 16.778 & -16.002 & -0.381 & -0.382\\
\midrule
08-05 & 1 & 3.758 & -2.287 & -0.061 & -0.061\\
 & 2 & 7.604 & -38.537 & -1.026 & -1.028\\
 & 4 & 15.451 & -61.531 & -1.603 & -1.603\\
\midrule
08-21 & 1 & 4.438 & 0.000 & 0.000 & 0.000\\
 & 2 & 9.011 & -1.024 & -0.023 & -0.023\\
 & 4 & 17.481 & -1.273 & -0.029 & -0.029\\
\bottomrule
\end{tabular*}
\vspace{3pt}
\parbox{\columnwidth}{Three-seed means. $E_0$ is recorded target-site energy without the signal over the corresponding window. $\delta P_C$ is the C-minus-R discrepancy in window-average power response. $\epsilon_m=100(\Delta E_m-\Delta E_R)/E_0$; ratios are computed within each pair before averaging. R, C, and A are defined in Section III.6.}
\end{table}
